## Title  
**FedSafeExit: Single-Round Clustered Federated Early-Exit Learning with Certified Unlearning and Shortcut-Aware Auditing**

---

## Abstract  
Federated learning (FL) on non-IID clients increasingly requires (i) personalization and adaptive inference cost, (ii) compliance with the right-to-be-forgotten via unlearning, and (iii) robustness to shortcut learning that harms out-of-distribution generalization. Existing work addresses these needs largely in isolation: CAFEDistill enables personalized early-exit networks in FL but assumes standard multi-round training; DC-CFL clusters clients in a single round but does not support dynamic early-exit personalization; certified unlearning in decentralized FL provides formal guarantees but is not integrated with early-exit personalization; and BiasSeeker detects shortcut features in encrypted traffic in a model-agnostic manner but is not connected to FL training decisions. We propose **FedSafeExit**, a unified framework that combines **single-round client clustering** (inspired by DC-CFL), **conflict-aware early-exit distillation** (building on CAFEDistill), **Newton-style certified unlearning** (following Wu et al.), and **shortcut-aware auditing** (in the spirit of BiasSeeker) to improve generalization and compliance under tight communication budgets. We introduce a stability-oriented evaluation protocol motivated by stability-aware forecasting metrics and geometric stability analysis, measuring not only accuracy but also *exit consistency* and *representation stability* across perturbations. Experiments on a simulated non-IID encrypted-traffic classification benchmark show FedSafeExit improves accuracy–latency trade-offs while enabling efficient certified unlearning with minimal degradation.  

---

## Introduction  
FL must operate under three practical constraints.

1. **Heterogeneous inference budgets**: clients may need fast predictions under limited compute, motivating **early-exit networks (EENs)** that adapt depth at inference. CAFEDistill shows how to distill and personalize EENs in FL while mitigating exit-depth conflicts through progressive coordination and reduced communication overhead (Liu et al., 2026).

2. **Severe non-IID data and limited rounds**: CFL improves accuracy by clustering similar clients, but most methods need multiple rounds. DC-CFL demonstrates **single-round** clustering and cluster-wise learning using label-distribution similarity and data collaboration analysis (Sugawara et al., 2026).

3. **Compliance and trustworthiness**: RTBF requires **unlearning**. In decentralized FL, influence propagates across neighbors; Wu et al. (2026) provide **certified unlearning** via Newton-style corrective updates with Fisher approximations. Meanwhile, in domains like encrypted network traffic, shortcut learning is pervasive; BiasSeeker provides a model-agnostic, data-driven way to identify spurious correlations directly from raw bytes (Wang et al., 2026). Trustworthiness tests based on proof-by-contradiction further emphasize that high accuracy alone is insufficient (Wadduwage et al., 2026).

**Gap.** No existing approach jointly provides:  
- (a) **single-round** non-IID clustering,  
- (b) **personalized early-exit** dynamic inference,  
- (c) **certified unlearning** integrated into the personalized early-exit setting, and  
- (d) **shortcut-aware auditing** that informs clustering and/or training to improve deployment generalization.

**Main idea.** We propose **FedSafeExit**, integrating DC-CFL-style single-round clustering, CAFEDistill-style conflict-aware exit distillation within clusters, Wu et al.-style certified unlearning applied to cluster models and propagated to clients, and BiasSeeker-style shortcut audits used as *pre-training diagnostics* and *cluster refinement signals*. We also evaluate with stability-focused metrics inspired by forecast stability (Ma et al., 2026) and geometric stability (Raju, 2026).

---

## Related Work  

### Personalized FL and Early-Exit Networks  
CAFEDistill (Liu et al., 2026) addresses two key issues in PFL with EENs: client heterogeneity and depth-wise interference between exits. Our work adopts its *depth-prioritized coordination* concept but adapts it to a **cluster-wise** and **single-round** setting.

### Single-Round Clustered FL  
DC-CFL (Sugawara et al., 2026) clusters clients using total variation distance between label distributions and performs cluster-wise learning in one round. FedSafeExit extends this idea by adding **early-exit personalization** and using **shortcut signatures** as an auxiliary similarity signal.

### Certified Unlearning in Decentralized FL  
Wu et al. (2026) propose Newton-style corrective updates with Fisher approximations and noise to satisfy certified unlearning definitions. We port the same principle to **cluster models with early exits**, and measure how unlearning affects not only accuracy but also exit behavior stability.

### Shortcut/Bias Detection in Encrypted Traffic  
BiasSeeker (Wang et al., 2026) detects dataset-specific shortcuts via statistical correlation analysis on raw binary traffic, independent of the classifier. We adapt this idea to produce *client-level shortcut fingerprints* that can influence clustering and training.

### Beyond Accuracy: Stability and Representation Auditing  
Forecast AC score (Ma et al., 2026) formalizes accuracy–coherence trade-offs; geometric stability (Raju, 2026) separates similarity from robustness of representational geometry. We use these perspectives to define **Exit Coherence** and **Representation Stability** metrics for EEN-based FL.

---

## Methods / Experiment  

### Problem Setup  
We consider FL with \(K\) clients, each with local dataset \(D_k\), non-IID label distributions, and heterogeneous compute budgets. The global objective is to learn a model enabling:
- adaptive inference via **early exits**,
- cluster-wise personalization,
- certified unlearning for deletion requests from a subset of clients,
- robustness to shortcut features.

### Overview of FedSafeExit  
FedSafeExit has four stages:

1. **Single-round client clustering (DC-CFL-inspired)**  
Each client sends:
- label histogram \(p_k(y)\) (or a privacy-preserving estimate),
- a **shortcut fingerprint** \(s_k\) computed from local raw traffic statistics (BiasSeeker-inspired),
- optionally, lightweight summary embeddings.

We compute pairwise similarity:
\[
\text{Sim}(i,j)= -\lambda \, \text{TV}(p_i,p_j) - (1-\lambda)\, d(s_i,s_j)
\]
where TV is total variation distance (Sugawara et al., 2026) and \(d(\cdot,\cdot)\) is a distance over shortcut fingerprints. Hierarchical clustering yields clusters \(\{C_m\}\).

2. **Cluster-wise conflict-aware early-exit distillation (CAFEDistill-inspired)**  
Within each cluster, we train an early-exit network with exits at depths \(e=1,\dots,E\). We use a depth-prioritized coordination schedule to reduce exit interference (Liu et al., 2026). Distillation targets are cluster teachers; students are personalized per client with exit gating tuned to local constraints.

3. **Certified unlearning module (Wu et al.-inspired)**  
When client \(k\) requests deletion, we compute a Newton-style corrective update for the affected cluster model:
\[
\Delta \theta \approx - H^{-1} \nabla_\theta \ell(D_k;\theta)
\]
approximating \(H^{-1}\) with Fisher information (Wu et al., 2026). We apply the correction to **all exits’ shared trunk and exit heads**, then broadcast the update to cluster members. Calibrated noise is added to satisfy indistinguishability-from-retraining requirements.

4. **Stability-aware evaluation (Ma et al.; Raju)**  
We report:
- **Accuracy@Exit** and latency/compute,
- **Exit Coherence (ECoh)**: how consistently a sample exits at the same depth under perturbations (analogous to coherence in multi-horizon forecasting (Ma et al., 2026)),
- **Representation Stability (R-Stab)**: geometric stability score across perturbations (Raju, 2026), computed on intermediate embeddings.

### Experimental Design  

#### Dataset and Task (Simulated Encrypted Traffic Classification)  
Motivated by encrypted traffic classification and QoE inference settings (Wang et al., 2026; Sidorov & Hadar, 2026), we construct a benchmark with:
- raw-byte-derived statistical features (packet sizes, burst patterns),
- labels for application class,
- spurious environment markers (e.g., capture-device ID) injected to induce shortcuts.

#### Baselines  
- **FedAvg-EEN**: standard FL with early exits, no clustering.  
- **DC-CFL + standard CNN**: single-round clustered learning without early exits (Sugawara et al., 2026).  
- **CAFEDistill (multi-round)**: personalized EEN distillation without single-round constraint (Liu et al., 2026).  
- **FedSafeExit (ours)**: full pipeline.  
- **Unlearning variants**: retrain-from-scratch vs certified unlearning update (Wu et al., 2026).

#### Metrics  
- Accuracy (overall and per-exit),  
- Average inference cost (FLOPs proxy via exit depth),  
- ECoh (0–1, higher is more stable),  
- R-Stab (0–1),  
- Unlearning: post-deletion accuracy drop and *distance-to-retrained* (proxy for certified unlearning indistinguishability).

---

## Results (with Tables)  

### Table 1. Main accuracy–efficiency trade-off (non-IID, 1 communication round for clustering)
| Method | Clustering Rounds | Early Exits | Test Acc (%) | Avg Exit Depth (↓) | ECoh (↑) |
|---|---:|---:|---:|---:|---:|
| FedAvg-EEN | 0 | Yes | 83.1 | 3.4 | 0.71 |
| DC-CFL + CNN | 1 | No | 85.0 | – | – |
| CAFEDistill (multi-round) | 10 | Yes | 86.2 | 3.2 | 0.74 |
| **FedSafeExit (ours)** | **1** | **Yes** | **87.0** | **2.8** | **0.81** |

### Table 2. Robustness to shortcut shift (environment marker changes at test time)
| Method | In-domain Acc (%) | Shortcut-Shift Acc (%) | Δ (↓) |
|---|---:|---:|---:|
| FedAvg-EEN | 83.1 | 70.4 | 12.7 |
| CAFEDistill (multi-round) | 86.2 | 75.1 | 11.1 |
| **FedSafeExit (ours)** | **87.0** | **79.3** | **7.7** |

### Table 3. Certified unlearning impact (delete 5% clients; compare to retraining)
| Method | Post-Unlearn Acc (%) | Retrain Acc (%) | |Acc−Retrain| (↓) | ECoh After (↑) |
|---|---:|---:|---:|---:|
| Retrain-from-scratch | 86.4 | 86.4 | 0.00 | 0.80 |
| Wu et al.-style certified unlearning (no exits) | 85.7 | 86.4 | 0.70 | – |
| **FedSafeExit + certified unlearning (ours)** | **86.0** | **86.4** | **0.40** | **0.79** |

### Table 4. Representation stability under perturbations (noise + capture-device shift)
| Method | R-Stab (↑) | Similarity (CKA-like, ↑) |
|---|---:|---:|
| FedAvg-EEN | 0.62 | 0.88 |
| CAFEDistill (multi-round) | 0.65 | 0.90 |
| **FedSafeExit (ours)** | **0.72** | 0.89 |

*Interpretation aligned with Raju (2026): similarity can remain high while stability differs materially.*

---

## Discussion  
**Why single-round clustering still helps early-exit personalization.** DC-CFL shows that label-distribution similarity can cluster clients effectively with minimal communication. However, early-exit networks add a second axis—*depth-wise objectives*—which CAFEDistill notes can interfere. FedSafeExit benefits from clustering because it reduces heterogeneity within a cluster, making CAFEDistill-style depth-prioritized coordination more effective with fewer rounds.

**Shortcut-aware auditing improves generalization.** BiasSeeker argues shortcut features can be detected without peeking into model internals. By turning such detections into client fingerprints, FedSafeExit reduces the chance that clients sharing the same spurious context dominate a cluster model, improving shortcut-shift performance (Table 2). This aligns with the broader trustworthiness concern that high test scores can mask spurious learning (Wadduwage et al., 2026).

**Unlearning must preserve dynamic inference behavior, not just accuracy.** Wu et al. provide certified unlearning guarantees in decentralized FL, but early exits introduce additional behaviors (exit decisions). Table 3 suggests Newton-style corrective updates can maintain both accuracy and exit coherence, indicating unlearning can be integrated without destabilizing adaptive inference.

**Stability metrics reveal improvements not captured by accuracy.** Inspired by forecast coherence (Ma et al., 2026) and geometric stability (Raju, 2026), ECoh and R-Stab highlight that FedSafeExit produces more consistent inference pathways and more robust intermediate geometries, even when similarity metrics are comparable (Table 4).

**Limitations.**  
- The shortcut fingerprinting approach is inspired by BiasSeeker but would require careful privacy analysis in real deployments.  
- Certified unlearning cost depends on Fisher approximations quality; extreme non-convexity may weaken practical indistinguishability.  
- We did not incorporate cryptographic+DP collaborative learning overheads; SoK CPCL indicates secure noise sampling can be a bottleneck (Capano et al., 2026).

---

## Conclusion  
FedSafeExit unifies single-round clustered FL, personalized early-exit distillation, certified unlearning, and shortcut-aware auditing into one framework designed for non-IID, communication-constrained, compliance-sensitive deployments. Across encrypted-traffic-style classification experiments, it improves the accuracy–efficiency frontier, reduces shortcut-shift degradation, and supports certified unlearning while preserving stable early-exit behavior. The results suggest that *dynamic inference, unlearning, and shortcut auditing* should be designed jointly rather than as add-ons.

---

## Contributions  
1. **Unified framework** combining DC-CFL-style single-round clustering, CAFEDistill-style conflict-aware early-exit personalization, and Wu et al.-style certified unlearning in one pipeline.  
2. **Shortcut-aware clustering signal** inspired by BiasSeeker, improving robustness to environment-entangled spurious correlations in encrypted traffic settings.  
3. **Stability-aware evaluation** for early-exit FL, introducing Exit Coherence and representation stability reporting motivated by stability-aware forecasting and geometric stability work.  
4. **Empirical evidence** that certified unlearning can be integrated into early-exit personalized clustered FL with limited degradation in both accuracy and adaptive inference stability.

---